\section{Probability Flow ODE}

\subsection{从 SDE 到 ODE}

前向 SDE 描述了一个随机过程，其轨迹是锯齿状的（因为 Wiener 过程的路径虽然连续但处处不可微）。一个自然的问题是：

\textit{是否存在一个确定性 ODE，使得它产生的边际分布 $p_t(\mathbf{x})$ 与 SDE 完全相同？}

\begin{importantbox}{Probability Flow ODE}
给定前向 SDE：
$$
\mathrm{d}\mathbf{x} = f(\mathbf{x}, t)\,\mathrm{d}t + g(t)\,\mathrm{d}\mathbf{w}
$$
存在唯一的 ODE（称为 \textbf{Probability Flow ODE}）：
$$
\mathrm{d}\mathbf{x} = \left[f(\mathbf{x}, t) - \frac{1}{2}g(t)^2 \nabla_{\mathbf{x}} \log p_t(\mathbf{x})\right]\mathrm{d}t
$$
使得该 ODE 的轨迹集合在每个时刻 $t$ 的边际分布与 SDE 的边际分布 $p_t(\mathbf{x})$ 一致。
\end{importantbox}

注意 Probability Flow ODE 与反向 SDE 的关系：

\begin{table}[H]
\centering
\begin{tabular}{lcc}
\toprule
& \textbf{漂移项修正} & \textbf{随机项} \\
\midrule
反向 SDE & $f - g^2 \nabla_{\mathbf{x}} \log p_t$ & $g(t)\,\mathrm{d}\bar{\mathbf{w}}$ \\
Probability Flow ODE & $f - \frac{1}{2} g^2 \nabla_{\mathbf{x}} \log p_t$ & 无 \\
\bottomrule
\end{tabular}
\caption{反向 SDE 与 Probability Flow ODE 的对比}
\end{table}

ODE 形式中 score 的系数从 $g^2$ 变为 $\frac{1}{2}g^2$，同时\textbf{去掉了 Wiener 过程项}。

\subsection{ODE 的优势}

Probability Flow ODE 带来了多项实际优势：

\begin{enumerate}
    \item \textbf{确定性采样}：给定初始噪声 $\mathbf{x}_T$，ODE 的解是确定的——同一个噪声总是生成同一张图像。这与 DDIM 的效果一致（事实上 DDIM 可以看作 Probability Flow ODE 的一种离散化）。
    \item \textbf{更少的采样步数}：确定性 ODE 的轨迹平滑，允许使用更大的步长和更高阶的数值求解器（如 Runge-Kutta 方法），从而显著减少采样步数。
    \item \textbf{精确的似然计算}：利用 instantaneous change of variables 公式，可以通过 ODE 精确计算数据的对数似然 $\log p(\mathbf{x}_0)$。
    \item \textbf{语义插值}：由于 ODE 建立了数据空间和噪声空间之间的双射（bijection），可以在潜空间中插值来获得语义上有意义的中间结果。
\end{enumerate}

\begin{knowledgebox}{DDIM 与 Probability Flow ODE 的联系}
DDIM 可以被理解为 VP-SDE 对应 Probability Flow ODE 的一种特殊离散化方案。这解释了 DDIM 的两个核心特性：（1）采样过程是确定性的（没有随机项），（2）可以使用远少于训练时的步数（如 50 步代替 1000 步）进行高质量生成。
\end{knowledgebox}

\subsection{ODE 轨迹 vs SDE 轨迹}

\begin{itemize}
    \item \textbf{SDE 轨迹}：锯齿状、随机，同一起点出发的不同轨迹到达不同终点。边际分布通过对所有可能轨迹取平均得到。
    \item \textbf{ODE 轨迹}：光滑、确定性，每个起点唯一对应一个终点。边际分布由轨迹簇的整体变形（probability flow）给出。
\end{itemize}

两者的边际分布 $p_t(\mathbf{x})$ 在每个时刻 $t$ 完全一致，但\textbf{单条轨迹}是不同的。

\begin{warningbox}{Probability Flow ODE 仍然需要 Score}
尽管 ODE 去掉了随机项，但 score function $\nabla_{\mathbf{x}} \log p_t(\mathbf{x})$ 仍然出现在漂移项中。因此，训练过程不变——无论最终使用 SDE sampler 还是 ODE sampler 进行采样，模型训练都是学习 score function（或等价地，学习噪声预测）。
\end{warningbox}

\subsection{与 Neural ODE 和 Normalizing Flows 的联系}

Probability Flow ODE 揭示了扩散模型与 Neural ODE（Chen et al., 2018）以及 Continuous Normalizing Flows（CNF）之间的深层联系。

\begin{knowledgebox}{Continuous Normalizing Flows}
CNF 通过一个由神经网络参数化的向量场 $v_\theta(\mathbf{x}, t)$ 定义 ODE $\mathrm{d}\mathbf{x}/\mathrm{d}t = v_\theta(\mathbf{x}, t)$，将简单分布（如高斯）映射到复杂数据分布。Probability Flow ODE 正是这种形式——其向量场为 $f(\mathbf{x}, t) - \frac{1}{2}g(t)^2 \nabla_{\mathbf{x}} \log p_t(\mathbf{x})$，其中 score 部分由神经网络近似。

关键区别在于：CNF 直接训练向量场，而扩散模型先通过 score matching 训练 score 网络，然后利用 Probability Flow ODE 进行采样。后者的训练更加稳定，因为 score matching 有明确的去噪目标。
\end{knowledgebox}

\begin{importantbox}{Instantaneous Change of Variables}
对于由 ODE $\mathrm{d}\mathbf{x}/\mathrm{d}t = \mathbf{v}(\mathbf{x}, t)$ 生成的流，对数密度的演化满足：
$$
\frac{\mathrm{d}\log p_t(\mathbf{x}_t)}{\mathrm{d}t} = -\mathrm{tr}\left(\frac{\partial \mathbf{v}}{\partial \mathbf{x}}\right)
$$
将 Probability Flow ODE 的向量场代入此公式，就可以精确计算数据的对数似然 $\log p_0(\mathbf{x}_0)$。这是扩散模型相比 GAN 的重要优势之一——GAN 无法提供精确的似然值。
\end{importantbox}

\subsection{本章小结}

Probability Flow ODE 是 SDE 的确定性"镜像"：它保留了相同的边际分布，但轨迹光滑、无随机性。这为确定性采样、高效求解和精确似然计算打开了大门，也建立了扩散模型与 Neural ODE / Continuous Normalizing Flows 之间的深层联系。


